% Ngân hàng bài mẫu: Phân thức đại số (lớp 8).
% Quy ước:
%   - Mỗi bài: \begin{problem}[id] ... \end{problem}; lời giải lồng bên trong.
%   - id: duy nhất, không dấu, không đổi, không dùng lại. Dạng <lớp><phân môn>-<chủ đề>-<số>.
%   - Mỗi bước lời giải là 1 đoạn, các bước cách nhau 1 dòng trống; mỗi bước có lý do.
%   - Dòng cuối lời giải: Đáp số: $\boxed{...}$
%   - Cách giải khác: thêm \begin{solution}[cách 2].
%   - Metadata bằng comment ngay sau \begin{problem}: % nguon:, % lop:
\documentclass[12pt,a4paper]{article}
\usepackage[utf8]{vietnam}
\usepackage{amsmath,amssymb}
\newenvironment{problem}[1][]{\par\medskip\noindent\textbf{Bài #1.}\ }{\par}
\newenvironment{solution}[1][]{\par\noindent\textit{Lời giải #1.}\par}{\par}

\begin{document}

% ---------------------------------------------------------------- điều kiện xác định
\begin{problem}[8DS-PS-001]
% nguon: dummy (tự sinh để test)
% lop: 8
% kc_du_kien: 8DS-PS-01, 8DS-HDT-02
% do_kho_du_kien: 1
Tìm điều kiện xác định của mỗi phân thức sau:

a) $\dfrac{3x}{x - 2}$;

b) $\dfrac{x + 1}{x^2 - 9}$.

\begin{solution}
a) Phân thức xác định khi mẫu thức khác $0$: $x - 2 \neq 0$.

a) Chuyển vế $-2$ sang vế phải: $x \neq 2$.

b) Phân thức xác định khi mẫu thức khác $0$: $x^2 - 9 \neq 0$.

b) Phân tích mẫu thành nhân tử theo hằng đẳng thức hiệu hai bình phương: $x^2 - 9 = (x - 3)(x + 3)$.

b) Tích $(x - 3)(x + 3)$ khác $0$ khi và chỉ khi cả hai thừa số đều khác $0$: $x - 3 \neq 0$ và $x + 3 \neq 0$, tức là $x \neq 3$ và $x \neq -3$.

Đáp số: a) $\boxed{x \neq 2}$; b) $\boxed{x \neq 3 \text{ và } x \neq -3}$
\end{solution}
\end{problem}

% ---------------------------------------------------------------- rút gọn phân thức
\begin{problem}[8DS-PS-002]
% nguon: dummy (tự sinh để test)
% lop: 8
% kc_du_kien: 8DS-PS-02, 8DS-PTNT-01, 8DS-PTNT-02
% do_kho_du_kien: 2
Rút gọn phân thức $\dfrac{3x^2 - 6x}{x^2 - 4x + 4}$.

\begin{solution}
Điều kiện xác định: mẫu $x^2 - 4x + 4 \neq 0$. Theo hằng đẳng thức bình phương của một hiệu, $x^2 - 4x + 4 = (x - 2)^2$, nên điều kiện là $x \neq 2$.

Phân tích tử thức thành nhân tử bằng cách đặt nhân tử chung $3x$: $3x^2 - 6x = 3x(x - 2)$.

Viết lại phân thức dưới dạng tích: $\dfrac{3x^2 - 6x}{x^2 - 4x + 4} = \dfrac{3x(x - 2)}{(x - 2)^2}$.

Chia cả tử và mẫu cho nhân tử chung $x - 2$ (khác $0$ do $x \neq 2$), theo tính chất cơ bản của phân thức: $\dfrac{3x(x - 2)}{(x - 2)^2} = \dfrac{3x}{x - 2}$.

Đáp số: $\boxed{\dfrac{3x}{x - 2} \text{ với } x \neq 2}$
\end{solution}
\end{problem}

% ---------------------------------------------------------------- cộng hai phân thức, đổi dấu, chữ
\begin{problem}[8DS-PS-003]
% nguon: dummy (tự sinh để test)
% lop: 8
% kc_du_kien: 8DS-PS-03, 8DS-PS-02
% do_kho_du_kien: 2
Thực hiện phép tính $\dfrac{x^2}{x - y} + \dfrac{y^2}{y - x}$.

\begin{solution}
Điều kiện xác định: $x - y \neq 0$ và $y - x \neq 0$, tức là $x \neq y$.

Đổi dấu cả tử và mẫu của phân thức thứ hai (quy tắc đổi dấu): $\dfrac{y^2}{y - x} = \dfrac{-y^2}{x - y}$.

Hai phân thức có cùng mẫu $x - y$, nên cộng các tử và giữ nguyên mẫu: $\dfrac{x^2}{x - y} + \dfrac{-y^2}{x - y} = \dfrac{x^2 - y^2}{x - y}$.

Phân tích tử theo hằng đẳng thức hiệu hai bình phương: $x^2 - y^2 = (x - y)(x + y)$.

Chia cả tử và mẫu cho $x - y$ (khác $0$ do $x \neq y$): $\dfrac{(x - y)(x + y)}{x - y} = x + y$.

Đáp số: $\boxed{x + y \text{ với } x \neq y}$
\end{solution}
\end{problem}

% ---------------------------------------------------------------- nhân và chia phân thức
\begin{problem}[8DS-PS-004]
% nguon: dummy (tự sinh để test)
% lop: 8
% kc_du_kien: 8DS-PS-04, 8DS-PS-02, 8DS-PS-01
% do_kho_du_kien: 3
Thực hiện các phép tính sau:

a) $\dfrac{x^2 - 9}{2x + 4} \cdot \dfrac{x + 2}{x^2 - 3x}$;

b) $\dfrac{x^2 - 4x + 4}{x + 1} : \dfrac{x - 2}{x^2 + x}$.

\begin{solution}
a) Phân tích các tử và mẫu thành nhân tử: $x^2 - 9 = (x - 3)(x + 3)$ (hiệu hai bình phương), $2x + 4 = 2(x + 2)$ và $x^2 - 3x = x(x - 3)$ (đặt nhân tử chung).

a) Điều kiện xác định: các mẫu khác $0$, tức là $2(x + 2) \neq 0$ và $x(x - 3) \neq 0$, hay $x \neq -2$, $x \neq 0$, $x \neq 3$.

a) Nhân tử với tử, mẫu với mẫu (quy tắc nhân phân thức):
\[ \dfrac{x^2 - 9}{2x + 4} \cdot \dfrac{x + 2}{x^2 - 3x} = \dfrac{(x - 3)(x + 3)(x + 2)}{2(x + 2) \cdot x(x - 3)}. \]

a) Chia cả tử và mẫu cho nhân tử chung $(x - 3)(x + 2)$ (khác $0$ theo điều kiện), được $\dfrac{x + 3}{2x}$.

b) Phân tích thành nhân tử: $x^2 - 4x + 4 = (x - 2)^2$ (bình phương của một hiệu) và $x^2 + x = x(x + 1)$ (đặt nhân tử chung).

b) Điều kiện xác định: các mẫu khác $0$ nên $x + 1 \neq 0$ và $x(x + 1) \neq 0$; ngoài ra phân thức chia $\dfrac{x - 2}{x^2 + x}$ phải khác $0$ nên $x - 2 \neq 0$. Vậy $x \neq -1$, $x \neq 0$, $x \neq 2$.

b) Chia cho một phân thức là nhân với phân thức nghịch đảo của nó:
\[ \dfrac{(x - 2)^2}{x + 1} : \dfrac{x - 2}{x(x + 1)} = \dfrac{(x - 2)^2}{x + 1} \cdot \dfrac{x(x + 1)}{x - 2} = \dfrac{(x - 2)^2 \cdot x(x + 1)}{(x + 1)(x - 2)}. \]

b) Chia cả tử và mẫu cho nhân tử chung $(x + 1)(x - 2)$ (khác $0$ theo điều kiện), được $x(x - 2)$.

Đáp số: a) $\boxed{\dfrac{x + 3}{2x} \text{ với } x \neq 0,\ x \neq -2,\ x \neq 3}$; b) $\boxed{x(x - 2) \text{ với } x \neq 0,\ x \neq -1,\ x \neq 2}$
\end{solution}
\end{problem}

% ---------------------------------------------------------------- cộng trừ ba phân thức
\begin{problem}[8DS-PS-005]
% nguon: dummy (tự sinh để test)
% lop: 8
% kc_du_kien: 8DS-PS-03, 8DS-PS-02
% do_kho_du_kien: 3
Thực hiện phép tính $\dfrac{x + 1}{2x - 2} - \dfrac{x - 1}{2x + 2} - \dfrac{2}{x^2 - 1}$.

\begin{solution}
Phân tích các mẫu thành nhân tử: $2x - 2 = 2(x - 1)$, $2x + 2 = 2(x + 1)$ (đặt nhân tử chung) và $x^2 - 1 = (x - 1)(x + 1)$ (hiệu hai bình phương).

Điều kiện xác định: các mẫu khác $0$, tức là $x \neq 1$ và $x \neq -1$.

Mẫu thức chung là $2(x - 1)(x + 1)$; nhân tử phụ của ba phân thức lần lượt là $x + 1$, $x - 1$ và $2$.

Quy đồng mẫu rồi trừ các tử, giữ nguyên mẫu chung:
\[ \dfrac{x + 1}{2x - 2} - \dfrac{x - 1}{2x + 2} - \dfrac{2}{x^2 - 1} = \dfrac{(x + 1)^2 - (x - 1)^2 - 4}{2(x - 1)(x + 1)}. \]

Khai triển tử theo hằng đẳng thức bình phương của một tổng, một hiệu: $(x^2 + 2x + 1) - (x^2 - 2x + 1) - 4 = 4x - 4 = 4(x - 1)$.

Chia cả tử và mẫu cho $2(x - 1)$ (khác $0$ vì $x \neq 1$): $\dfrac{4(x - 1)}{2(x - 1)(x + 1)} = \dfrac{2}{x + 1}$.

Đáp số: $\boxed{\dfrac{2}{x + 1} \text{ với } x \neq 1,\ x \neq -1}$
\end{solution}
\end{problem}

% ---------------------------------------------------------------- bài toán thực tế, có đơn vị
\begin{problem}[8DS-PS-006]
% nguon: dummy (tự sinh để test)
% lop: 8
% kc_du_kien: 8DS-PS-03, 8DS-PS-05
% do_kho_du_kien: 3
Một ô tô đi từ A đến B trên quãng đường dài $120$ km với vận tốc $x$ (km/h). Lúc về, ô tô tăng vận tốc thêm $10$ km/h.

a) Viết biểu thức biểu thị tổng thời gian cả đi và về của ô tô theo $x$ dưới dạng một phân thức.

b) Tính tổng thời gian đó khi $x = 30$.

\begin{solution}
a) Vận tốc là số dương nên $x > 0$; khi đó các mẫu $x$ và $x + 10$ đều khác $0$.

a) Thời gian bằng quãng đường chia cho vận tốc: thời gian đi là $\dfrac{120}{x}$ (giờ); vận tốc lúc về là $x + 10$ (km/h) nên thời gian về là $\dfrac{120}{x + 10}$ (giờ).

a) Tổng thời gian là $T = \dfrac{120}{x} + \dfrac{120}{x + 10}$. Quy đồng với mẫu thức chung $x(x + 10)$ rồi cộng các tử:
\[ T = \dfrac{120(x + 10) + 120x}{x(x + 10)} = \dfrac{240x + 1200}{x(x + 10)} = \dfrac{240(x + 5)}{x(x + 10)} \text{ (giờ)}. \]

b) Giá trị $x = 30$ thoả điều kiện $x > 0$. Thay vào phân thức vừa tìm: $T = \dfrac{240 \cdot (30 + 5)}{30 \cdot (30 + 10)} = \dfrac{8\,400}{1\,200} = 7$ (giờ).

b) Kiểm tra lại: thời gian đi $\dfrac{120}{30} = 4$ giờ, thời gian về $\dfrac{120}{40} = 3$ giờ, tổng $7$ giờ (đúng).

Đáp số: a) $\boxed{T = \dfrac{240(x + 5)}{x(x + 10)} \text{ (giờ)}}$; b) $\boxed{7 \text{ giờ}}$
\end{solution}
\end{problem}

% ---------------------------------------------------------------- rút gọn rồi tính giá trị, đối chiếu ĐKXĐ
\begin{problem}[8DS-PS-007]
% nguon: dummy (tự sinh để test)
% lop: 8
% kc_du_kien: 8DS-PS-05, 8DS-PS-01, 8DS-PTNT-05
% do_kho_du_kien: 4
Cho biểu thức $P = \left(\dfrac{1}{x + 1} + \dfrac{2}{x^2 - 1}\right) \cdot \dfrac{x^2 - 2x + 1}{x}$.

a) Tìm điều kiện xác định và rút gọn $P$.

b) Tính giá trị của $P$ biết $x^2 - 4x + 3 = 0$.

\begin{solution}
a) Các mẫu thức là $x + 1$, $x^2 - 1 = (x - 1)(x + 1)$ và $x$. Điều kiện xác định: các mẫu khác $0$, tức là $x \neq 0$, $x \neq 1$, $x \neq -1$.

a) Quy đồng trong ngoặc với mẫu thức chung $(x - 1)(x + 1)$, rồi chia tử và mẫu cho $x + 1 \neq 0$:
\[ \dfrac{1}{x + 1} + \dfrac{2}{(x - 1)(x + 1)} = \dfrac{(x - 1) + 2}{(x - 1)(x + 1)} = \dfrac{x + 1}{(x - 1)(x + 1)} = \dfrac{1}{x - 1}. \]

a) Theo hằng đẳng thức bình phương của một hiệu, $x^2 - 2x + 1 = (x - 1)^2$. Nhân hai phân thức rồi chia tử và mẫu cho $x - 1 \neq 0$: $P = \dfrac{1}{x - 1} \cdot \dfrac{(x - 1)^2}{x} = \dfrac{(x - 1)^2}{(x - 1)x} = \dfrac{x - 1}{x}$.

b) Phân tích vế trái thành nhân tử bằng cách tách $-4x = -x - 3x$: $x^2 - 4x + 3 = x^2 - x - 3x + 3 = x(x - 1) - 3(x - 1) = (x - 1)(x - 3)$.

b) Tích $(x - 1)(x - 3)$ bằng $0$ khi một trong hai thừa số bằng $0$: $x - 1 = 0$ hoặc $x - 3 = 0$, tức là $x = 1$ hoặc $x = 3$.

b) Đối chiếu điều kiện xác định: $x = 1$ không thoả mãn điều kiện $x \neq 1$ nên loại; $x = 3$ thoả mãn.

b) Thay $x = 3$ vào biểu thức đã rút gọn: $P = \dfrac{3 - 1}{3} = \dfrac{2}{3}$.

Đáp số: a) $\boxed{P = \dfrac{x - 1}{x} \text{ với } x \neq 0,\ x \neq 1,\ x \neq -1}$; b) $\boxed{P = \dfrac{2}{3}}$
\end{solution}
\end{problem}

% ---------------------------------------------------------------- chứng minh không phụ thuộc biến
\begin{problem}[8DS-PS-008]
% nguon: dummy (tự sinh để test)
% lop: 8
% kc_du_kien: 8DS-PS-05, 8DS-PS-03, 8DS-PS-04
% do_kho_du_kien: 4
Chứng minh rằng với $x \neq 3$ và $x \neq -3$, giá trị của biểu thức
\[ Q = \left(\dfrac{2x}{x + 3} + \dfrac{x}{x - 3} - \dfrac{3x^2 + 9}{x^2 - 9}\right) : \dfrac{3}{3 - x} \]
không phụ thuộc vào giá trị của biến $x$.

\begin{solution}
Với $x \neq 3$ và $x \neq -3$, các mẫu $x + 3$, $x - 3$, $x^2 - 9 = (x - 3)(x + 3)$ và $3 - x$ đều khác $0$; phân thức chia $\dfrac{3}{3 - x}$ có tử $3 \neq 0$ nên khác $0$. Do đó biểu thức $Q$ xác định.

Quy đồng trong ngoặc với mẫu thức chung $(x - 3)(x + 3)$; nhân tử phụ lần lượt là $x - 3$, $x + 3$ và $1$:
\[ \dfrac{2x}{x + 3} + \dfrac{x}{x - 3} - \dfrac{3x^2 + 9}{x^2 - 9} = \dfrac{2x(x - 3) + x(x + 3) - (3x^2 + 9)}{(x - 3)(x + 3)}. \]

Bỏ ngoặc và thu gọn tử: $2x^2 - 6x + x^2 + 3x - 3x^2 - 9 = -3x - 9 = -3(x + 3)$.

Chia cả tử và mẫu cho $x + 3 \neq 0$: biểu thức trong ngoặc bằng $\dfrac{-3(x + 3)}{(x - 3)(x + 3)} = \dfrac{-3}{x - 3}$.

Chia cho một phân thức là nhân với phân thức nghịch đảo: $Q = \dfrac{-3}{x - 3} \cdot \dfrac{3 - x}{3} = \dfrac{-3(3 - x)}{3(x - 3)}$.

Vì $-(3 - x) = x - 3$ nên $Q = \dfrac{3(x - 3)}{3(x - 3)} = 1$ (chia tử và mẫu cho $3(x - 3) \neq 0$).

Vậy với mọi $x \neq 3$, $x \neq -3$ ta có $Q = 1$, giá trị của $Q$ không phụ thuộc vào $x$.

Đáp số: $\boxed{Q = 1}$
\end{solution}
\end{problem}

% ---------------------------------------------------------------- có cách giải khác: tách phân thức
\begin{problem}[8DS-PS-009]
% nguon: dummy (tự sinh để test)
% lop: 8
% kc_du_kien: 8DS-PS-03, 8DS-PS-02
% do_kho_du_kien: 4
Rút gọn biểu thức
\[ M = \dfrac{1}{x(x + 1)} + \dfrac{1}{(x + 1)(x + 2)} + \dfrac{1}{(x + 2)(x + 3)}. \]

\begin{solution}[cách 1: quy đồng mẫu]
Điều kiện xác định: các mẫu khác $0$, tức là $x \neq 0$, $x \neq -1$, $x \neq -2$, $x \neq -3$.

Mẫu thức chung là $x(x + 1)(x + 2)(x + 3)$; nhân tử phụ của ba phân thức lần lượt là $(x + 2)(x + 3)$, $x(x + 3)$ và $x(x + 1)$.

Quy đồng rồi cộng các tử, giữ nguyên mẫu chung: $M = \dfrac{(x + 2)(x + 3) + x(x + 3) + x(x + 1)}{x(x + 1)(x + 2)(x + 3)}$.

Bỏ ngoặc và thu gọn tử: $(x^2 + 5x + 6) + (x^2 + 3x) + (x^2 + x) = 3x^2 + 9x + 6 = 3(x^2 + 3x + 2)$.

Phân tích $x^2 + 3x + 2$ thành nhân tử bằng cách tách $3x = x + 2x$: $x^2 + x + 2x + 2 = x(x + 1) + 2(x + 1) = (x + 1)(x + 2)$. Vậy tử bằng $3(x + 1)(x + 2)$.

Chia cả tử và mẫu cho $(x + 1)(x + 2) \neq 0$: $M = \dfrac{3(x + 1)(x + 2)}{x(x + 1)(x + 2)(x + 3)} = \dfrac{3}{x(x + 3)}$.

Đáp số: $\boxed{M = \dfrac{3}{x(x + 3)} \text{ với } x \neq 0,\ x \neq -1,\ x \neq -2,\ x \neq -3}$
\end{solution}

\begin{solution}[cách 2: tách mỗi phân thức thành hiệu]
Điều kiện xác định: các mẫu khác $0$, tức là $x \neq 0$, $x \neq -1$, $x \neq -2$, $x \neq -3$.

Với $k \neq 0$ và $k \neq -1$, quy đồng hai phân thức ta có $\dfrac{1}{k} - \dfrac{1}{k + 1} = \dfrac{(k + 1) - k}{k(k + 1)} = \dfrac{1}{k(k + 1)}$.

Áp dụng đẳng thức trên với $k = x$, $k = x + 1$, $k = x + 2$:
\[ M = \left(\dfrac{1}{x} - \dfrac{1}{x + 1}\right) + \left(\dfrac{1}{x + 1} - \dfrac{1}{x + 2}\right) + \left(\dfrac{1}{x + 2} - \dfrac{1}{x + 3}\right). \]

Các cặp phân thức đối nhau có tổng bằng $0$ nên triệt tiêu: $M = \dfrac{1}{x} - \dfrac{1}{x + 3}$.

Quy đồng với mẫu thức chung $x(x + 3)$ rồi trừ các tử: $M = \dfrac{(x + 3) - x}{x(x + 3)} = \dfrac{3}{x(x + 3)}$.

Đáp số: $\boxed{M = \dfrac{3}{x(x + 3)} \text{ với } x \neq 0,\ x \neq -1,\ x \neq -2,\ x \neq -3}$
\end{solution}
\end{problem}

% ---------------------------------------------------------------- tổng hợp: rút gọn, tính giá trị, giá trị nguyên
\begin{problem}[8DS-PS-010]
% nguon: dummy (tự sinh để test)
% lop: 8
% kc_du_kien: 8DS-PS-05, 8DS-PS-01, 8DS-PS-03
% do_kho_du_kien: 5
Cho biểu thức
\[ A = \left(\dfrac{x}{x^2 - 4} + \dfrac{2}{2 - x} + \dfrac{1}{x + 2}\right) : \left(x - 2 + \dfrac{10 - x^2}{x + 2}\right). \]
a) Tìm điều kiện xác định và rút gọn $A$.

b) Tính giá trị của $A$ tại $x$ thoả mãn $|x| = \dfrac{1}{2}$.

c) Tìm các giá trị nguyên của $x$ để $A$ có giá trị nguyên.

\begin{solution}
a) Các mẫu là $x^2 - 4 = (x - 2)(x + 2)$, $2 - x$ và $x + 2$; chúng khác $0$ khi $x \neq 2$ và $x \neq -2$.

a) Biến đổi ngoặc thứ hai bằng cách quy đồng với mẫu $x + 2$: $x - 2 + \dfrac{10 - x^2}{x + 2} = \dfrac{(x - 2)(x + 2) + 10 - x^2}{x + 2} = \dfrac{x^2 - 4 + 10 - x^2}{x + 2} = \dfrac{6}{x + 2}$. Phân thức này có tử $6 \neq 0$ nên luôn khác $0$; vậy điều kiện xác định của $A$ là $x \neq 2$ và $x \neq -2$.

a) Biến đổi ngoặc thứ nhất: đổi dấu $\dfrac{2}{2 - x} = \dfrac{-2}{x - 2}$, rồi quy đồng với mẫu thức chung $(x - 2)(x + 2)$:
\[ \dfrac{x}{(x - 2)(x + 2)} - \dfrac{2}{x - 2} + \dfrac{1}{x + 2} = \dfrac{x - 2(x + 2) + (x - 2)}{(x - 2)(x + 2)} = \dfrac{-6}{(x - 2)(x + 2)}. \]

a) Chia là nhân với phân thức nghịch đảo, rồi chia tử và mẫu cho $6(x + 2) \neq 0$: $A = \dfrac{-6}{(x - 2)(x + 2)} \cdot \dfrac{x + 2}{6} = \dfrac{-1}{x - 2} = \dfrac{1}{2 - x}$.

b) Từ $|x| = \dfrac{1}{2}$ suy ra $x = \dfrac{1}{2}$ hoặc $x = -\dfrac{1}{2}$ (định nghĩa giá trị tuyệt đối); cả hai giá trị đều thoả điều kiện $x \neq 2$, $x \neq -2$.

b) Với $x = \dfrac{1}{2}$, thay vào biểu thức đã rút gọn: $A = \dfrac{1}{2 - \frac{1}{2}} = \dfrac{1}{\frac{3}{2}} = \dfrac{2}{3}$.

b) Với $x = -\dfrac{1}{2}$: $A = \dfrac{1}{2 + \frac{1}{2}} = \dfrac{1}{\frac{5}{2}} = \dfrac{2}{5}$.

c) Với $x$ nguyên, $x \neq 2$ thì $2 - x$ là số nguyên khác $0$. Khi đó $A = \dfrac{1}{2 - x}$ là số nguyên khi và chỉ khi $2 - x$ là ước của $1$, tức là $2 - x = 1$ hoặc $2 - x = -1$.

c) Giải: $2 - x = 1$ cho $x = 1$; $2 - x = -1$ cho $x = 3$.

c) Đối chiếu điều kiện: $x = 1$ và $x = 3$ đều khác $2$ và $-2$ (thoả mãn). Thử lại: $x = 1$ cho $A = 1$; $x = 3$ cho $A = -1$ (đều là số nguyên).

Đáp số: a) $\boxed{A = \dfrac{1}{2 - x} \text{ với } x \neq 2,\ x \neq -2}$; b) $\boxed{A = \dfrac{2}{3} \text{ khi } x = \dfrac{1}{2};\ A = \dfrac{2}{5} \text{ khi } x = -\dfrac{1}{2}}$; c) $\boxed{x = 1 \text{ hoặc } x = 3}$
\end{solution}
\end{problem}

% ---------------------------------------------------------------- giá trị của phân thức tại một giá trị của biến
\begin{problem}[8DS-PS-011]
% nguon: dummy (tự sinh để test)
% lop: 8
% kc_du_kien: 8DS-PS-01, 8DS-PS-05
% do_kho_du_kien: 1
Tính giá trị của phân thức $P = \dfrac{x^2 - 4}{x + 5}$ tại $x = -1$ và tại $x = -5$ (nếu có).

\begin{solution}
Điều kiện xác định: mẫu thức khác $0$, tức là $x + 5 \neq 0$, hay $x \neq -5$.

Giá trị $x = -1$ thoả mãn điều kiện xác định. Thay $x = -1$ vào phân thức: $P = \dfrac{(-1)^2 - 4}{-1 + 5} = \dfrac{1 - 4}{4} = -\dfrac{3}{4}$.

Giá trị $x = -5$ không thoả mãn điều kiện xác định (mẫu bằng $-5 + 5 = 0$), nên tại $x = -5$ phân thức $P$ không xác định và không có giá trị.

Đáp số: $\boxed{P = -\dfrac{3}{4} \text{ tại } x = -1;\ P \text{ không xác định tại } x = -5}$
\end{solution}
\end{problem}

% ---------------------------------------------------------------- hai phân thức bằng nhau, tìm đa thức chưa biết
\begin{problem}[8DS-PS-012]
% nguon: dummy (tự sinh để test)
% lop: 8
% kc_du_kien: 8DS-PS-06, 8DS-PTNT-02, 8DS-PTNT-01
% do_kho_du_kien: 2
a) Dùng định nghĩa hai phân thức bằng nhau, chứng tỏ rằng $\dfrac{x - 3}{x + 3} = \dfrac{x^2 - 6x + 9}{x^2 - 9}$ với $x \neq 3$, $x \neq -3$.

b) Tìm đa thức $A$ thoả mãn $\dfrac{A}{2x - 6} = \dfrac{x^2 + 3x}{x^2 - 9}$ với $x \neq 3$, $x \neq -3$.

\begin{solution}
a) Với $x \neq 3$, $x \neq -3$, các mẫu $x + 3$ và $x^2 - 9 = (x - 3)(x + 3)$ đều khác $0$. Theo định nghĩa, $\dfrac{A}{B} = \dfrac{C}{D}$ khi $A \cdot D = B \cdot C$, nên ta so sánh hai tích chéo.

a) Tích thứ nhất, phân tích $x^2 - 9$ theo hằng đẳng thức hiệu hai bình phương: $(x - 3)(x^2 - 9) = (x - 3)(x - 3)(x + 3) = (x - 3)^2(x + 3)$.

a) Tích thứ hai, viết $x^2 - 6x + 9$ theo hằng đẳng thức bình phương của một hiệu: $(x + 3)(x^2 - 6x + 9) = (x + 3)(x - 3)^2$.

a) Hai tích chéo bằng nhau nên theo định nghĩa, hai phân thức đã cho bằng nhau.

b) Theo định nghĩa hai phân thức bằng nhau: $A \cdot (x^2 - 9) = (2x - 6)(x^2 + 3x)$.

b) Phân tích vế phải thành nhân tử (đặt nhân tử chung, rồi dùng hiệu hai bình phương): $(2x - 6)(x^2 + 3x) = 2(x - 3) \cdot x(x + 3) = 2x(x - 3)(x + 3) = 2x(x^2 - 9)$.

b) Vậy $A \cdot (x^2 - 9) = 2x \cdot (x^2 - 9)$. Hai vế là tích của cùng đa thức $x^2 - 9$ (khác $0$) với $A$ và với $2x$, nên $A = 2x$.

b) Thử lại: $\dfrac{2x}{2x - 6} = \dfrac{2x}{2(x - 3)} = \dfrac{x}{x - 3}$ và $\dfrac{x^2 + 3x}{x^2 - 9} = \dfrac{x(x + 3)}{(x - 3)(x + 3)} = \dfrac{x}{x - 3}$ (đúng).

Đáp số: a) $\boxed{(x - 3)(x^2 - 9) = (x + 3)(x^2 - 6x + 9) \text{ nên hai phân thức bằng nhau}}$; b) $\boxed{A = 2x}$
\end{solution}
\end{problem}

% ---------------------------------------------------------------- tìm x để phân thức bằng 0, đối chiếu ĐKXĐ
\begin{problem}[8DS-PS-013]
% nguon: dummy (tự sinh để test)
% lop: 8
% kc_du_kien: 8DS-PS-01, 8DS-PTNT-06, 8DS-PTNT-01
% do_kho_du_kien: 2
Tìm giá trị của $x$ để giá trị của phân thức $\dfrac{x^2 - 4x}{x^2 - 16}$ bằng $0$.

\begin{solution}
Điều kiện xác định: $x^2 - 16 \neq 0$. Theo hằng đẳng thức hiệu hai bình phương, $x^2 - 16 = (x - 4)(x + 4)$, nên điều kiện là $x \neq 4$ và $x \neq -4$.

Một phân thức xác định có giá trị bằng $0$ khi và chỉ khi tử thức bằng $0$: $x^2 - 4x = 0$.

Đặt nhân tử chung $x$: $x(x - 4) = 0$. Tích bằng $0$ khi một trong hai thừa số bằng $0$: $x = 0$ hoặc $x = 4$.

Đối chiếu điều kiện xác định: $x = 4$ không thoả mãn ($x \neq 4$) nên loại; $x = 0$ thoả mãn.

Thử lại: với $x = 0$, phân thức bằng $\dfrac{0}{-16} = 0$ (đúng).

Đáp số: $\boxed{x = 0}$
\end{solution}
\end{problem}

% ---------------------------------------------------------------- quy đồng mẫu thức nhiều phân thức
\begin{problem}[8DS-PS-014]
% nguon: dummy (tự sinh để test)
% lop: 8
% kc_du_kien: 8DS-PS-03, 8DS-PTNT-01, 8DS-PTNT-02
% do_kho_du_kien: 2
Quy đồng mẫu thức các phân thức $\dfrac{1}{x^2 + 2x}$, $\dfrac{3}{x^2 - 4}$ và $\dfrac{x}{2x - 4}$.

\begin{solution}
Phân tích các mẫu thành nhân tử: $x^2 + 2x = x(x + 2)$ và $2x - 4 = 2(x - 2)$ (đặt nhân tử chung); $x^2 - 4 = (x - 2)(x + 2)$ (hiệu hai bình phương).

Điều kiện xác định: các mẫu khác $0$, tức là $x \neq 0$, $x \neq 2$, $x \neq -2$.

Mẫu thức chung là tích của các nhân tử $2$, $x$, $x - 2$, $x + 2$, mỗi nhân tử lấy với số mũ lớn nhất: $2x(x - 2)(x + 2)$.

Nhân tử phụ của mỗi mẫu bằng mẫu thức chung chia cho mẫu đó: của $x(x + 2)$ là $2(x - 2)$; của $(x - 2)(x + 2)$ là $2x$; của $2(x - 2)$ là $x(x + 2)$.

Nhân cả tử và mẫu của mỗi phân thức với nhân tử phụ tương ứng (tính chất cơ bản của phân thức):
\[ \dfrac{1}{x^2 + 2x} = \dfrac{2(x - 2)}{2x(x - 2)(x + 2)}, \qquad \dfrac{3}{x^2 - 4} = \dfrac{6x}{2x(x - 2)(x + 2)}, \qquad \dfrac{x}{2x - 4} = \dfrac{x^2(x + 2)}{2x(x - 2)(x + 2)}. \]

Đáp số: $\boxed{\dfrac{2(x - 2)}{2x(x - 2)(x + 2)},\ \dfrac{6x}{2x(x - 2)(x + 2)},\ \dfrac{x^2(x + 2)}{2x(x - 2)(x + 2)}}$
\end{solution}
\end{problem}

% ---------------------------------------------------------------- rút gọn phân thức hai biến, hiệu hai lập phương
\begin{problem}[8DS-PS-015]
% nguon: dummy (tự sinh để test)
% lop: 8
% kc_du_kien: 8DS-PS-02, 8DS-PTNT-05, 8DS-HDT-04
% do_kho_du_kien: 3
Rút gọn phân thức $\dfrac{x^2 + xy - 2y^2}{x^3 - y^3}$.

\begin{solution}
Phân tích mẫu theo hằng đẳng thức hiệu hai lập phương: $x^3 - y^3 = (x - y)(x^2 + xy + y^2)$.

Điều kiện xác định: $x^3 - y^3 \neq 0$, tức là $x^3 \neq y^3$; hai số có lập phương bằng nhau thì bằng nhau, nên điều kiện là $x \neq y$. Khi đó tích $(x - y)(x^2 + xy + y^2)$ khác $0$ nên cả hai thừa số đều khác $0$.

Phân tích tử thành nhân tử bằng cách tách $xy = 2xy - xy$ rồi nhóm hạng tử: $x^2 + 2xy - xy - 2y^2 = x(x + 2y) - y(x + 2y) = (x - y)(x + 2y)$.

Chia cả tử và mẫu cho nhân tử chung $x - y$ (khác $0$ do $x \neq y$), theo tính chất cơ bản của phân thức:
\[ \dfrac{(x - y)(x + 2y)}{(x - y)(x^2 + xy + y^2)} = \dfrac{x + 2y}{x^2 + xy + y^2}. \]

Đáp số: $\boxed{\dfrac{x + 2y}{x^2 + xy + y^2} \text{ với } x \neq y}$
\end{solution}
\end{problem}

% ---------------------------------------------------------------- bài toán thực tế: tỉ số thời gian (chia phân thức)
\begin{problem}[8DS-PS-016]
% nguon: dummy (tự sinh để test)
% lop: 8
% kc_du_kien: 8DS-PS-04, 8DS-PS-05
% do_kho_du_kien: 3
Hai tổ sản xuất, mỗi tổ phải làm $360$ sản phẩm. Mỗi giờ tổ I làm được $x$ sản phẩm, tổ II làm được ít hơn tổ I $6$ sản phẩm ($x > 6$).

a) Viết các phân thức biểu thị thời gian hoàn thành công việc của mỗi tổ.

b) Viết phân thức (đã rút gọn) biểu thị tỉ số giữa thời gian hoàn thành của tổ II và thời gian hoàn thành của tổ I.

c) Tính tỉ số đó khi $x = 30$.

\begin{solution}
a) Thời gian hoàn thành bằng số sản phẩm chia cho số sản phẩm làm được mỗi giờ. Tổ I làm $x$ sản phẩm mỗi giờ nên thời gian là $\dfrac{360}{x}$ (giờ).

a) Tổ II mỗi giờ làm $x - 6$ sản phẩm ($x - 6 > 0$ do $x > 6$) nên thời gian là $\dfrac{360}{x - 6}$ (giờ). Vì $x > 6$ nên các mẫu $x$ và $x - 6$ đều khác $0$.

b) Tỉ số cần tìm là thương $\dfrac{360}{x - 6} : \dfrac{360}{x}$. Chia cho một phân thức là nhân với phân thức nghịch đảo của nó: $\dfrac{360}{x - 6} \cdot \dfrac{x}{360} = \dfrac{360x}{360(x - 6)}$.

b) Chia cả tử và mẫu cho $360$: tỉ số bằng $\dfrac{x}{x - 6}$.

c) Giá trị $x = 30$ thoả mãn $x > 6$. Thay vào phân thức đã rút gọn: $\dfrac{30}{30 - 6} = \dfrac{30}{24} = \dfrac{5}{4}$.

c) Kiểm tra: tổ I làm trong $360 : 30 = 12$ giờ, tổ II làm trong $360 : 24 = 15$ giờ, và $\dfrac{15}{12} = \dfrac{5}{4}$ (đúng).

Đáp số: a) $\boxed{\dfrac{360}{x} \text{ (giờ) và } \dfrac{360}{x - 6} \text{ (giờ)}}$; b) $\boxed{\dfrac{x}{x - 6}}$; c) $\boxed{\dfrac{5}{4}}$
\end{solution}
\end{problem}

% ---------------------------------------------------------------- chứng minh đẳng thức hai biến
\begin{problem}[8DS-PS-017]
% nguon: dummy (tự sinh để test)
% lop: 8
% kc_du_kien: 8DS-PS-03, 8DS-PS-04, 8DS-HDT-02
% do_kho_du_kien: 3
Chứng minh rằng với $x \neq 0$, $y \neq 0$, $x \neq y$ và $x \neq -y$ thì
\[ \left(\dfrac{1}{x} - \dfrac{1}{y}\right) : \left(\dfrac{1}{x^2} - \dfrac{1}{y^2}\right) = \dfrac{xy}{x + y}. \]

\begin{solution}
Quy đồng ngoặc thứ nhất với mẫu thức chung $xy$ (khác $0$ do $x \neq 0$, $y \neq 0$): $\dfrac{1}{x} - \dfrac{1}{y} = \dfrac{y - x}{xy}$.

Quy đồng ngoặc thứ hai với mẫu thức chung $x^2y^2$, rồi phân tích tử theo hằng đẳng thức hiệu hai bình phương: $\dfrac{1}{x^2} - \dfrac{1}{y^2} = \dfrac{y^2 - x^2}{x^2y^2} = \dfrac{(y - x)(y + x)}{x^2y^2}$.

Phân thức chia khác $0$ vì tử $(y - x)(y + x) \neq 0$ (do $x \neq y$ và $x \neq -y$), nên phép chia thực hiện được.

Chia cho một phân thức là nhân với phân thức nghịch đảo của nó:
\[ \dfrac{y - x}{xy} : \dfrac{(y - x)(y + x)}{x^2y^2} = \dfrac{y - x}{xy} \cdot \dfrac{x^2y^2}{(y - x)(y + x)} = \dfrac{(y - x) \cdot x^2y^2}{xy(y - x)(y + x)}. \]

Chia cả tử và mẫu cho $xy(y - x)$ (khác $0$ theo giả thiết): vế trái bằng $\dfrac{xy}{y + x} = \dfrac{xy}{x + y}$, đúng bằng vế phải. Đẳng thức được chứng minh.

Đáp số: $\boxed{\left(\dfrac{1}{x} - \dfrac{1}{y}\right) : \left(\dfrac{1}{x^2} - \dfrac{1}{y^2}\right) = \dfrac{xy}{x + y}}$
\end{solution}
\end{problem}

% ---------------------------------------------------------------- tính giá trị biểu thức khi biến thoả điều kiện, có cách giải khác
\begin{problem}[8DS-PS-018]
% nguon: dummy (tự sinh để test)
% lop: 8
% kc_du_kien: 8DS-PS-05, 8DS-PS-02
% do_kho_du_kien: 4
Cho các số $x$, $y$, $z$ thoả mãn $3x - y = 3z$, $2x + y = 7z$ và $z \neq 0$. Tính giá trị của biểu thức $M = \dfrac{x^2 - 2xy}{x^2 + y^2}$.

\begin{solution}[cách 1: biểu diễn $x$, $y$ theo $z$]
Cộng vế theo vế hai đẳng thức đã cho: $(3x - y) + (2x + y) = 3z + 7z$, tức là $5x = 10z$; chia hai vế cho $5$ được $x = 2z$.

Thay $x = 2z$ vào đẳng thức $3x - y = 3z$: $6z - y = 3z$, suy ra $y = 3z$.

Kiểm tra điều kiện xác định: mẫu $x^2 + y^2 = 4z^2 + 9z^2 = 13z^2 \neq 0$ vì $z \neq 0$, nên $M$ xác định.

Tính tử: $x^2 - 2xy = 4z^2 - 2 \cdot 2z \cdot 3z = 4z^2 - 12z^2 = -8z^2$.

Chia cả tử và mẫu cho $z^2 \neq 0$: $M = \dfrac{-8z^2}{13z^2} = -\dfrac{8}{13}$.

Đáp số: $\boxed{M = -\dfrac{8}{13}}$
\end{solution}

\begin{solution}[cách 2: khử $z$, tính theo tỉ số $\dfrac{x}{y}$]
Hai đẳng thức cho $21z = 7(3x - y)$ và $21z = 3(2x + y)$, nên $7(3x - y) = 3(2x + y)$. Bỏ ngoặc: $21x - 7y = 6x + 3y$; chuyển vế: $15x = 10y$, tức là $3x = 2y$.

Chứng tỏ $y \neq 0$: nếu $y = 0$ thì từ $3x = 2y$ có $x = 0$, khi đó $3z = 3x - y = 0$ nên $z = 0$, trái giả thiết. Vậy $y \neq 0$ và $\dfrac{x}{y} = \dfrac{2}{3}$; mẫu $x^2 + y^2 \geq y^2 > 0$ nên $M$ xác định.

Chia cả tử và mẫu của $M$ cho $y^2 \neq 0$ (tính chất cơ bản của phân thức), đặt $t = \dfrac{x}{y}$: $M = \dfrac{\frac{x^2}{y^2} - \frac{2x}{y}}{\frac{x^2}{y^2} + 1} = \dfrac{t^2 - 2t}{t^2 + 1}$.

Thay $t = \dfrac{2}{3}$: tử $\dfrac{4}{9} - \dfrac{4}{3} = -\dfrac{8}{9}$, mẫu $\dfrac{4}{9} + 1 = \dfrac{13}{9}$, nên $M = -\dfrac{8}{9} : \dfrac{13}{9} = -\dfrac{8}{13}$.

Đáp số: $\boxed{M = -\dfrac{8}{13}}$
\end{solution}
\end{problem}

% ---------------------------------------------------------------- giá trị lớn nhất của biểu thức hữu tỉ
\begin{problem}[8DS-PS-019]
% nguon: dummy (tự sinh để test)
% lop: 8
% kc_du_kien: 8DS-PS-05, 8DS-PS-01, 8DS-HDT-01
% do_kho_du_kien: 4
Tìm giá trị lớn nhất của biểu thức $B = \dfrac{x^2 - 4x + 10}{x^2 - 4x + 7}$.

\begin{solution}
Biến đổi mẫu theo hằng đẳng thức bình phương của một hiệu: $x^2 - 4x + 7 = (x^2 - 4x + 4) + 3 = (x - 2)^2 + 3$.

Vì $(x - 2)^2 \geq 0$ với mọi $x$ nên mẫu $(x - 2)^2 + 3 \geq 3 > 0$; do đó $B$ xác định với mọi $x$.

Tách tử: $x^2 - 4x + 10 = (x^2 - 4x + 7) + 3$, nên $B = \dfrac{(x^2 - 4x + 7) + 3}{x^2 - 4x + 7} = 1 + \dfrac{3}{(x - 2)^2 + 3}$.

Hai phân thức có cùng tử dương, phân thức nào có mẫu dương lớn hơn thì nhỏ hơn. Từ $(x - 2)^2 + 3 \geq 3$ suy ra $\dfrac{3}{(x - 2)^2 + 3} \leq \dfrac{3}{3} = 1$, nên $B \leq 1 + 1 = 2$ với mọi $x$.

Dấu bằng xảy ra khi $(x - 2)^2 = 0$, tức là $x = 2$. Thử lại: $B = \dfrac{4 - 8 + 10}{4 - 8 + 7} = \dfrac{6}{3} = 2$ (đúng).

Đáp số: $\boxed{\max B = 2 \text{ khi } x = 2}$
\end{solution}
\end{problem}

% ---------------------------------------------------------------- biểu thức ba biến có điều kiện a + b + c = 0
\begin{problem}[8DS-PS-020]
% nguon: dummy (tự sinh để test)
% lop: 8
% kc_du_kien: 8DS-PS-05, 8DS-PS-03, 8DS-HDT-01
% do_kho_du_kien: 5
Cho ba số $a$, $b$, $c$ đều khác $0$ và thoả mãn $a + b + c = 0$. Tính giá trị của biểu thức
\[ P = \dfrac{1}{b^2 + c^2 - a^2} + \dfrac{1}{c^2 + a^2 - b^2} + \dfrac{1}{a^2 + b^2 - c^2}. \]

\begin{solution}
Từ $a + b + c = 0$ suy ra $a = -(b + c)$, nên theo hằng đẳng thức bình phương của một tổng $a^2 = (b + c)^2 = b^2 + 2bc + c^2$. Do đó $b^2 + c^2 - a^2 = -2bc$.

Tương tự, $b = -(c + a)$ nên $b^2 = c^2 + 2ca + a^2$, do đó $c^2 + a^2 - b^2 = -2ca$.

Tương tự, $c = -(a + b)$ nên $c^2 = a^2 + 2ab + b^2$, do đó $a^2 + b^2 - c^2 = -2ab$.

Vì $a$, $b$, $c$ đều khác $0$ nên các mẫu $-2bc$, $-2ca$, $-2ab$ đều khác $0$; biểu thức $P$ xác định.

Thay vào: $P = -\dfrac{1}{2bc} - \dfrac{1}{2ca} - \dfrac{1}{2ab}$. Quy đồng với mẫu thức chung $2abc$ (nhân tử phụ lần lượt là $a$, $b$, $c$): $P = -\dfrac{a + b + c}{2abc}$.

Vì $a + b + c = 0$ và $2abc \neq 0$ nên $P = -\dfrac{0}{2abc} = 0$.

Đáp số: $\boxed{P = 0}$
\end{solution}
\end{problem}

% ---------------------------------------------------------------- nhân, chia phân thức có tử và mẫu là đơn thức
\begin{problem}[8DS-PS-021]
% nguon: dummy (tự sinh để test)
% lop: 8
% kc_du_kien: 8DS-PS-04, 8DS-PS-02
% do_kho_du_kien: 1
Thực hiện các phép tính sau:

a) $\dfrac{4x^2}{5y^3} \cdot \dfrac{15y}{8x}$;

b) $\dfrac{6x^3}{7y} : \dfrac{3x}{14y^2}$.

\begin{solution}
a) Điều kiện xác định: các mẫu $5y^3$ và $8x$ khác $0$, tức là $x \neq 0$ và $y \neq 0$.

a) Nhân tử với tử, mẫu với mẫu (quy tắc nhân phân thức): $\dfrac{4x^2}{5y^3} \cdot \dfrac{15y}{8x} = \dfrac{4x^2 \cdot 15y}{5y^3 \cdot 8x} = \dfrac{60x^2y}{40xy^3}$.

a) Chia cả tử và mẫu cho nhân tử chung $20xy$ (khác $0$ do $x \neq 0$, $y \neq 0$), theo tính chất cơ bản của phân thức: $\dfrac{60x^2y}{40xy^3} = \dfrac{3x}{2y^2}$.

b) Điều kiện xác định: các mẫu $7y$, $14y^2$ khác $0$ và phân thức chia $\dfrac{3x}{14y^2}$ khác $0$ (có tử $3x \neq 0$), tức là $x \neq 0$ và $y \neq 0$.

b) Chia cho một phân thức là nhân với phân thức nghịch đảo của nó: $\dfrac{6x^3}{7y} : \dfrac{3x}{14y^2} = \dfrac{6x^3}{7y} \cdot \dfrac{14y^2}{3x} = \dfrac{84x^3y^2}{21xy}$.

b) Chia cả tử và mẫu cho nhân tử chung $21xy$ (khác $0$ theo điều kiện): $\dfrac{84x^3y^2}{21xy} = 4x^2y$.

Đáp số: a) $\boxed{\dfrac{3x}{2y^2}}$; b) $\boxed{4x^2y}$
\end{solution}
\end{problem}

% ---------------------------------------------------------------- hai phân thức bằng nhau, tìm mẫu thức chưa biết
\begin{problem}[8DS-PS-022]
% nguon: dummy (tự sinh để test)
% lop: 8
% kc_du_kien: 8DS-PS-06, 8DS-HDT-04
% do_kho_du_kien: 2
Tìm đa thức $B$ thoả mãn $\dfrac{x^3 + 8}{B} = \dfrac{x^2 - 2x + 4}{x - 3}$ (với các giá trị của $x$ làm các mẫu khác $0$).

\begin{solution}
Theo định nghĩa hai phân thức bằng nhau ($\dfrac{A}{B} = \dfrac{C}{D}$ khi $A \cdot D = B \cdot C$), ta có $(x^3 + 8)(x - 3) = B \cdot (x^2 - 2x + 4)$.

Phân tích $x^3 + 8$ theo hằng đẳng thức tổng hai lập phương: $x^3 + 8 = x^3 + 2^3 = (x + 2)(x^2 - 2x + 4)$. Do đó $B \cdot (x^2 - 2x + 4) = (x + 2)(x - 3) \cdot (x^2 - 2x + 4)$.

Đa thức $x^2 - 2x + 4 = (x - 1)^2 + 3$ luôn dương nên khác $0$; chia hai vế cho $x^2 - 2x + 4$ được $B = (x + 2)(x - 3)$.

Nhân hai đa thức: $B = x^2 - 3x + 2x - 6 = x^2 - x - 6$.

Thử lại: với $x \neq -2$, $x \neq 3$ thì $\dfrac{x^3 + 8}{x^2 - x - 6} = \dfrac{(x + 2)(x^2 - 2x + 4)}{(x + 2)(x - 3)} = \dfrac{x^2 - 2x + 4}{x - 3}$ (đúng).

Đáp số: $\boxed{B = x^2 - x - 6}$
\end{solution}
\end{problem}

% ---------------------------------------------------------------- rút gọn phân thức nhiều chữ, nhóm hạng tử
\begin{problem}[8DS-PS-023]
% nguon: dummy (tự sinh để test)
% lop: 8
% kc_du_kien: 8DS-PS-02, 8DS-PTNT-03
% do_kho_du_kien: 2
Rút gọn phân thức $\dfrac{ax - ay + 2x - 2y}{ax + ay + 2x + 2y}$.

\begin{solution}
Phân tích tử thành nhân tử bằng cách nhóm hạng tử rồi đặt nhân tử chung: $ax - ay + 2x - 2y = a(x - y) + 2(x - y) = (x - y)(a + 2)$.

Phân tích mẫu tương tự: $ax + ay + 2x + 2y = a(x + y) + 2(x + y) = (x + y)(a + 2)$.

Điều kiện xác định: mẫu $(x + y)(a + 2) \neq 0$, tức là $x + y \neq 0$ và $a + 2 \neq 0$, hay $x \neq -y$ và $a \neq -2$.

Chia cả tử và mẫu cho nhân tử chung $a + 2$ (khác $0$ theo điều kiện), theo tính chất cơ bản của phân thức: $\dfrac{(x - y)(a + 2)}{(x + y)(a + 2)} = \dfrac{x - y}{x + y}$.

Đáp số: $\boxed{\dfrac{x - y}{x + y} \text{ với } x \neq -y,\ a \neq -2}$
\end{solution}
\end{problem}

% ---------------------------------------------------------------- hai phân thức bằng nhau, kiểm tra từng phân thức
\begin{problem}[8DS-PS-024]
% nguon: dummy (tự sinh để test)
% lop: 8
% kc_du_kien: 8DS-PS-06, 8DS-PTNT-01, 8DS-HDT-02
% do_kho_du_kien: 2
Dùng định nghĩa hai phân thức bằng nhau, cho biết mỗi phân thức sau có bằng phân thức $\dfrac{x - 2}{x + 2}$ hay không:

a) $\dfrac{2 - x}{-x - 2}$;

b) $\dfrac{x^2 - 2x}{x^2 + 2x}$;

c) $\dfrac{x^2 - 4}{x^2 + 4}$.

\begin{solution}
a) Theo định nghĩa, $\dfrac{A}{B} = \dfrac{C}{D}$ khi $A \cdot D = B \cdot C$. Tích chéo thứ nhất: $(x - 2)(-x - 2) = -(x - 2)(x + 2)$; tích chéo thứ hai: $(x + 2)(2 - x) = -(x + 2)(x - 2)$.

a) Hai tích chéo bằng nhau nên $\dfrac{2 - x}{-x - 2} = \dfrac{x - 2}{x + 2}$.

b) Đặt nhân tử chung: $x^2 + 2x = x(x + 2)$ và $x^2 - 2x = x(x - 2)$. Tích chéo thứ nhất: $(x - 2) \cdot x(x + 2) = x(x - 2)(x + 2)$; tích chéo thứ hai: $(x + 2) \cdot x(x - 2) = x(x - 2)(x + 2)$.

b) Hai tích chéo bằng nhau nên $\dfrac{x^2 - 2x}{x^2 + 2x} = \dfrac{x - 2}{x + 2}$.

c) Tích chéo thứ nhất: $(x - 2)(x^2 + 4) = x^3 - 2x^2 + 4x - 8$. Tích chéo thứ hai, dùng hằng đẳng thức hiệu hai bình phương $x^2 - 4 = (x - 2)(x + 2)$: $(x + 2)(x^2 - 4) = (x + 2)^2(x - 2) = x^3 + 2x^2 - 4x - 8$.

c) Hiệu hai tích chéo là $(x^3 + 2x^2 - 4x - 8) - (x^3 - 2x^2 + 4x - 8) = 4x^2 - 8x = 4x(x - 2)$, không phải đa thức $0$; chẳng hạn tại $x = 1$ tích thứ nhất bằng $-5$ còn tích thứ hai bằng $-9$. Vậy hai tích chéo không bằng nhau và hai phân thức không bằng nhau.

Đáp số: a) $\boxed{\dfrac{2 - x}{-x - 2} = \dfrac{x - 2}{x + 2}}$; b) $\boxed{\dfrac{x^2 - 2x}{x^2 + 2x} = \dfrac{x - 2}{x + 2}}$; c) $\boxed{\dfrac{x^2 - 4}{x^2 + 4} \neq \dfrac{x - 2}{x + 2}}$
\end{solution}
\end{problem}

% ---------------------------------------------------------------- tìm x để phân thức bằng một số, đối chiếu ĐKXĐ
\begin{problem}[8DS-PS-025]
% nguon: dummy (tự sinh để test)
% lop: 8
% kc_du_kien: 8DS-PS-01, 8DS-PS-02, 8DS-PT1-01
% do_kho_du_kien: 3
Cho phân thức $P = \dfrac{x^2 - 9}{x^2 - 3x}$.

a) Tìm giá trị của $x$ để $P = 2$.

b) Tìm giá trị của $x$ để $P = 4$.

\begin{solution}
a) Điều kiện xác định: mẫu $x^2 - 3x \neq 0$. Đặt nhân tử chung, $x^2 - 3x = x(x - 3)$, nên điều kiện là $x \neq 0$ và $x \neq 3$.

a) Rút gọn với điều kiện trên: theo hằng đẳng thức hiệu hai bình phương, $x^2 - 9 = (x - 3)(x + 3)$; chia cả tử và mẫu cho $x - 3 \neq 0$ được $P = \dfrac{(x - 3)(x + 3)}{x(x - 3)} = \dfrac{x + 3}{x}$.

a) $P = 2$ khi $\dfrac{x + 3}{x} = 2$. Nhân hai vế với $x \neq 0$: $x + 3 = 2x$; chuyển vế được $x = 3$.

a) Đối chiếu điều kiện xác định: $x = 3$ không thoả mãn $x \neq 3$ (tại đó mẫu $x^2 - 3x$ bằng $0$, phân thức $P$ không xác định) nên loại. Vậy không có giá trị nào của $x$ để $P = 2$.

b) Dùng phân thức đã rút gọn ở câu a: $P = 4$ khi $\dfrac{x + 3}{x} = 4$. Nhân hai vế với $x \neq 0$: $x + 3 = 4x$; chuyển vế được $3x = 3$, tức là $x = 1$.

b) Đối chiếu điều kiện xác định: $x = 1$ thoả mãn $x \neq 0$ và $x \neq 3$. Thử lại với phân thức ban đầu: $P = \dfrac{1 - 9}{1 - 3} = \dfrac{-8}{-2} = 4$ (đúng).

Đáp số: a) $\boxed{\text{không có giá trị nào của } x}$; b) $\boxed{x = 1}$
\end{solution}
\end{problem}

% ---------------------------------------------------------------- tìm phân thức chưa biết trong một tích
\begin{problem}[8DS-PS-026]
% nguon: dummy (tự sinh để test)
% lop: 8
% kc_du_kien: 8DS-PS-04, 8DS-PS-02, 8DS-PS-01
% do_kho_du_kien: 3
Tìm phân thức $Q$ (viết ở dạng rút gọn) thoả mãn
\[ \dfrac{x^2 + 2x}{x - 1} \cdot Q = \dfrac{x^2 - 4}{x^2 - 1}. \]

\begin{solution}
Phân tích các tử và mẫu thành nhân tử: $x^2 + 2x = x(x + 2)$ (đặt nhân tử chung); $x^2 - 4 = (x - 2)(x + 2)$ và $x^2 - 1 = (x - 1)(x + 1)$ (hiệu hai bình phương).

Điều kiện: các mẫu $x - 1$ và $x^2 - 1$ khác $0$ nên $x \neq 1$, $x \neq -1$; để tìm $Q$ phải chia cho phân thức $\dfrac{x^2 + 2x}{x - 1}$, nên phân thức này khác $0$, tức là $x(x + 2) \neq 0$, hay $x \neq 0$, $x \neq -2$.

$Q$ là một thừa số chưa biết của tích, nên $Q$ bằng tích chia cho thừa số đã biết: $Q = \dfrac{x^2 - 4}{x^2 - 1} : \dfrac{x^2 + 2x}{x - 1}$.

Chia cho một phân thức là nhân với phân thức nghịch đảo của nó, rồi nhân tử với tử, mẫu với mẫu:
\[ Q = \dfrac{(x - 2)(x + 2)}{(x - 1)(x + 1)} \cdot \dfrac{x - 1}{x(x + 2)} = \dfrac{(x - 2)(x + 2)(x - 1)}{(x - 1)(x + 1) \cdot x(x + 2)}. \]

Chia cả tử và mẫu cho nhân tử chung $(x - 1)(x + 2)$ (khác $0$ theo điều kiện): $Q = \dfrac{x - 2}{x(x + 1)}$.

Thử lại: $\dfrac{x(x + 2)}{x - 1} \cdot \dfrac{x - 2}{x(x + 1)} = \dfrac{(x + 2)(x - 2)}{(x - 1)(x + 1)} = \dfrac{x^2 - 4}{x^2 - 1}$ (đúng).

Đáp số: $\boxed{Q = \dfrac{x - 2}{x(x + 1)} \text{ với } x \neq 0,\ x \neq 1,\ x \neq -1,\ x \neq -2}$
\end{solution}
\end{problem}

% ---------------------------------------------------------------- bài toán thực tế: diện tích, số luống (nhân và chia phân thức)
\begin{problem}[8DS-PS-027]
% nguon: dummy (tự sinh để test)
% lop: 8
% kc_du_kien: 8DS-PS-04, 8DS-PS-05, 8DS-HDT-02
% do_kho_du_kien: 3
Một mảnh vườn hình chữ nhật có chiều dài $\dfrac{x^2 - 1}{x}$ (m) và chiều rộng $\dfrac{x^2}{x + 1}$ (m), với $x$ là số nguyên, $x > 2$.

a) Viết phân thức (đã rút gọn) biểu thị diện tích mảnh vườn theo $x$.

b) Người ta chia mảnh vườn thành các luống rau có cùng diện tích $\dfrac{x - 1}{2}$ (m$^2$). Viết biểu thức (đã rút gọn) biểu thị số luống rau theo $x$.

c) Tính diện tích mảnh vườn và số luống rau khi $x = 5$.

\begin{solution}
a) Vì $x > 2$ nên các mẫu $x$ và $x + 1$ đều dương, khác $0$; tử $x^2 - 1 > 0$ và $x^2 > 0$ nên chiều dài và chiều rộng đều là số dương.

a) Diện tích hình chữ nhật bằng chiều dài nhân chiều rộng. Phân tích $x^2 - 1 = (x - 1)(x + 1)$ (hiệu hai bình phương) rồi nhân tử với tử, mẫu với mẫu: $S = \dfrac{x^2 - 1}{x} \cdot \dfrac{x^2}{x + 1} = \dfrac{(x - 1)(x + 1) \cdot x^2}{x(x + 1)}$.

a) Chia cả tử và mẫu cho nhân tử chung $x(x + 1) \neq 0$: $S = x(x - 1)$ (m$^2$).

b) Số luống rau bằng diện tích mảnh vườn chia cho diện tích mỗi luống: $n = x(x - 1) : \dfrac{x - 1}{2}$; phép chia thực hiện được vì $x - 1 > 0$ nên phân thức chia khác $0$.

b) Chia cho một phân thức là nhân với phân thức nghịch đảo của nó, rồi chia tử và mẫu cho $x - 1 \neq 0$: $n = x(x - 1) \cdot \dfrac{2}{x - 1} = \dfrac{2x(x - 1)}{x - 1} = 2x$ (luống). Vì $x$ là số nguyên dương nên $2x$ là số tự nhiên (thoả mãn).

c) Giá trị $x = 5$ là số nguyên và $x > 2$ (thoả mãn). Thay vào các biểu thức đã rút gọn: $S = 5 \cdot (5 - 1) = 20$ (m$^2$) và $n = 2 \cdot 5 = 10$ (luống).

c) Kiểm tra: chiều dài $\dfrac{24}{5}$ m, chiều rộng $\dfrac{25}{6}$ m, diện tích $\dfrac{24}{5} \cdot \dfrac{25}{6} = 20$ (m$^2$); mỗi luống có diện tích $\dfrac{5 - 1}{2} = 2$ (m$^2$) và $20 : 2 = 10$ luống (đúng).

Đáp số: a) $\boxed{S = x(x - 1)\ (\text{m}^2)}$; b) $\boxed{2x \text{ luống}}$; c) $\boxed{S = 20\ \text{m}^2;\ 10 \text{ luống}}$
\end{solution}
\end{problem}

% ---------------------------------------------------------------- tham số để phân thức xác định với mọi x
\begin{problem}[8DS-PS-028]
% nguon: dummy (tự sinh để test)
% lop: 8
% kc_du_kien: 8DS-PS-01, 8DS-HDT-01
% do_kho_du_kien: 4
Tìm tất cả các giá trị của tham số $m$ để phân thức $\dfrac{x + 3}{x^2 + 2x + m}$ xác định với mọi giá trị của $x$.

\begin{solution}
Phân thức xác định khi mẫu thức khác $0$; vậy yêu cầu bài toán là $x^2 + 2x + m \neq 0$ với mọi giá trị của $x$.

Biến đổi mẫu theo hằng đẳng thức bình phương của một tổng: $x^2 + 2x + m = (x^2 + 2x + 1) + (m - 1) = (x + 1)^2 + m - 1$.

Trường hợp $m > 1$: vì $(x + 1)^2 \geq 0$ với mọi $x$ nên mẫu $(x + 1)^2 + m - 1 \geq m - 1 > 0$; mẫu luôn khác $0$, phân thức xác định với mọi $x$ (thoả mãn).

Trường hợp $m \leq 1$: số $1 - m \geq 0$ nên có căn bậc hai số học $\sqrt{1 - m}$. Chọn $x = -1 + \sqrt{1 - m}$ thì $(x + 1)^2 = 1 - m$, khi đó mẫu bằng $(1 - m) + m - 1 = 0$; phân thức không xác định tại giá trị này của $x$ (không thoả mãn). Chẳng hạn $m = 1$ thì mẫu bằng $(x + 1)^2$, bằng $0$ tại $x = -1$.

Kết hợp hai trường hợp: phân thức xác định với mọi $x$ khi và chỉ khi $m > 1$.

Đáp số: $\boxed{m > 1}$
\end{solution}
\end{problem}

% ---------------------------------------------------------------- rút gọn bằng nhóm hạng tử, thêm bớt hạng tử
\begin{problem}[8DS-PS-029]
% nguon: dummy (tự sinh để test)
% lop: 8
% kc_du_kien: 8DS-PS-02, 8DS-PTNT-04, 8DS-HDT-02
% do_kho_du_kien: 4
Rút gọn các phân thức sau:

a) $\dfrac{x^2 + y^2 - z^2 + 2xy}{x^2 - y^2 + z^2 + 2xz}$;

b) $\dfrac{x^4 + 4}{x^2 + 2x + 2}$.

\begin{solution}
a) Phân tích tử: nhóm $x^2 + 2xy + y^2$ thành bình phương của một tổng, rồi dùng hằng đẳng thức hiệu hai bình phương: $x^2 + 2xy + y^2 - z^2 = (x + y)^2 - z^2 = (x + y - z)(x + y + z)$.

a) Phân tích mẫu tương tự: $x^2 + 2xz + z^2 - y^2 = (x + z)^2 - y^2 = (x - y + z)(x + y + z)$.

a) Điều kiện xác định: mẫu $(x - y + z)(x + y + z) \neq 0$, tức là $x - y + z \neq 0$ và $x + y + z \neq 0$.

a) Chia cả tử và mẫu cho nhân tử chung $x + y + z$ (khác $0$ theo điều kiện), theo tính chất cơ bản của phân thức: $\dfrac{(x + y - z)(x + y + z)}{(x - y + z)(x + y + z)} = \dfrac{x + y - z}{x - y + z}$.

b) Mẫu: $x^2 + 2x + 2 = (x + 1)^2 + 1 \geq 1 > 0$ với mọi $x$ (bình phương của một tổng), nên phân thức xác định với mọi $x$.

b) Phân tích tử: thêm và bớt $4x^2$ để xuất hiện bình phương của một tổng: $x^4 + 4 = x^4 + 4x^2 + 4 - 4x^2 = (x^2 + 2)^2 - (2x)^2$.

b) Theo hằng đẳng thức hiệu hai bình phương: $(x^2 + 2)^2 - (2x)^2 = (x^2 + 2 - 2x)(x^2 + 2 + 2x) = (x^2 - 2x + 2)(x^2 + 2x + 2)$.

b) Chia cả tử và mẫu cho nhân tử chung $x^2 + 2x + 2$ (luôn khác $0$): $\dfrac{(x^2 - 2x + 2)(x^2 + 2x + 2)}{x^2 + 2x + 2} = x^2 - 2x + 2$.

Đáp số: a) $\boxed{\dfrac{x + y - z}{x - y + z} \text{ với } x + y + z \neq 0,\ x - y + z \neq 0}$; b) $\boxed{x^2 - 2x + 2}$
\end{solution}
\end{problem}

% ---------------------------------------------------------------- hai phân thức bằng nhau, tìm hệ số; có cách giải khác
\begin{problem}[8DS-PS-030]
% nguon: dummy (tự sinh để test)
% lop: 8
% kc_du_kien: 8DS-PS-06, 8DS-PS-03
% do_kho_du_kien: 4
Tìm các số $a$ và $b$ sao cho với mọi $x \neq 1$ và $x \neq -1$ ta có
\[ \dfrac{a}{x - 1} + \dfrac{b}{x + 1} = \dfrac{5x + 1}{x^2 - 1}. \]

\begin{solution}[cách 1: đồng nhất hệ số]
Quy đồng vế trái với mẫu thức chung $(x - 1)(x + 1) = x^2 - 1$ (nhân tử phụ lần lượt là $x + 1$ và $x - 1$): $\dfrac{a}{x - 1} + \dfrac{b}{x + 1} = \dfrac{a(x + 1) + b(x - 1)}{x^2 - 1} = \dfrac{(a + b)x + (a - b)}{x^2 - 1}$.

Hai phân thức $\dfrac{(a + b)x + (a - b)}{x^2 - 1}$ và $\dfrac{5x + 1}{x^2 - 1}$ có cùng mẫu khác $0$; theo định nghĩa hai phân thức bằng nhau, chúng bằng nhau khi $[(a + b)x + (a - b)](x^2 - 1) = (x^2 - 1)(5x + 1)$, tức là khi hai tử bằng nhau: $(a + b)x + (a - b) = 5x + 1$.

Hai đa thức bằng nhau khi các hệ số của các hạng tử cùng bậc bằng nhau: $a + b = 5$ và $a - b = 1$.

Cộng vế theo vế hai đẳng thức: $2a = 6$, nên $a = 3$; suy ra $b = 5 - 3 = 2$.

Thử lại: $\dfrac{3}{x - 1} + \dfrac{2}{x + 1} = \dfrac{3(x + 1) + 2(x - 1)}{x^2 - 1} = \dfrac{5x + 1}{x^2 - 1}$ (đúng).

Đáp số: $\boxed{a = 3,\ b = 2}$
\end{solution}

\begin{solution}[cách 2: thay giá trị cụ thể của $x$]
Quy đồng vế trái với mẫu thức chung $x^2 - 1$ như cách 1 rồi so sánh hai tử (định nghĩa hai phân thức bằng nhau): đẳng thức đã cho đúng với mọi $x \neq \pm 1$ khi $a(x + 1) + b(x - 1) = 5x + 1$ với mọi $x \neq \pm 1$.

Thay $x = 0$ (thoả mãn điều kiện): $a \cdot 1 + b \cdot (-1) = 1$, tức là $a - b = 1$.

Thay $x = 2$ (thoả mãn điều kiện): $a \cdot 3 + b \cdot 1 = 11$, tức là $3a + b = 11$.

Cộng vế theo vế hai đẳng thức vừa tìm: $4a = 12$, nên $a = 3$; suy ra $b = a - 1 = 2$.

Các bước trên mới cho điều kiện cần, nên phải kiểm tra lại: với $a = 3$, $b = 2$ thì $3(x + 1) + 2(x - 1) = 5x + 1$ đúng với mọi $x$, nên đẳng thức đã cho đúng với mọi $x \neq \pm 1$.

Đáp số: $\boxed{a = 3,\ b = 2}$
\end{solution}
\end{problem}

% ---------------------------------------------------------------- rút gọn phân thức ba biến, hằng đẳng thức bậc ba
\begin{problem}[8DS-PS-031]
% nguon: dummy (tự sinh để test)
% lop: 8
% kc_du_kien: 8DS-PS-02, 8DS-HDT-04, 8DS-HDT-03
% do_kho_du_kien: 5
Rút gọn phân thức
\[ \dfrac{a^3 + b^3 + c^3 - 3abc}{(a - b)^2 + (b - c)^2 + (c - a)^2}. \]

\begin{solution}
Điều kiện xác định: mẫu là tổng của ba bình phương nên không âm, và bằng $0$ khi và chỉ khi $a - b = b - c = c - a = 0$, tức là $a = b = c$. Vậy phân thức xác định khi $a$, $b$, $c$ không đồng thời bằng nhau.

Theo hằng đẳng thức lập phương của một tổng, $(a + b)^3 = a^3 + b^3 + 3ab(a + b)$, nên $a^3 + b^3 = (a + b)^3 - 3ab(a + b)$. Thay vào tử: $a^3 + b^3 + c^3 - 3abc = \left[(a + b)^3 + c^3\right] - 3ab(a + b + c)$.

Theo hằng đẳng thức tổng hai lập phương: $(a + b)^3 + c^3 = (a + b + c)\left[(a + b)^2 - (a + b)c + c^2\right] = (a + b + c)(a^2 + 2ab + b^2 - ac - bc + c^2)$.

Đặt nhân tử chung $a + b + c$: tử bằng $(a + b + c)(a^2 + 2ab + b^2 - ac - bc + c^2 - 3ab) = (a + b + c)(a^2 + b^2 + c^2 - ab - bc - ca)$.

Khai triển mẫu theo hằng đẳng thức bình phương của một hiệu: $(a^2 - 2ab + b^2) + (b^2 - 2bc + c^2) + (c^2 - 2ca + a^2) = 2(a^2 + b^2 + c^2 - ab - bc - ca)$.

Theo điều kiện xác định, mẫu khác $0$ nên $a^2 + b^2 + c^2 - ab - bc - ca \neq 0$. Chia cả tử và mẫu cho nhân tử chung này:
\[ \dfrac{(a + b + c)(a^2 + b^2 + c^2 - ab - bc - ca)}{2(a^2 + b^2 + c^2 - ab - bc - ca)} = \dfrac{a + b + c}{2}. \]

Đáp số: $\boxed{\dfrac{a + b + c}{2} \text{ với } a,\ b,\ c \text{ không đồng thời bằng nhau}}$
\end{solution}
\end{problem}

\end{document}
